% Paper-ready methods supplement for AdverSim.
% This file is intentionally package-light. It uses tables and verbatim-style
% pseudocode so it can be pasted into IEEE or ACM templates with minimal edits.

\section{Algorithmic Methodology and Complexity}

Let $N$ be the number of scenarios, $R$ the maximum number of corrective
re-query attempts, $L$ the generated JSON length, $S$ the number of stages in a
scenario, $T$ the number of ATT\&CK technique mentions, $G$ the number of attack
graph edges, $K$ the number of active ATT\&CK techniques, $P$ the number of
ATT\&CK group co-occurrence pairs, $M$ the number of Sigma rules, $E$ the number
of replay events, $H$ the number of human reviewers, and $B$ the number of
bootstrap or permutation resamples. LLM inference dominates wall-clock
generation cost; the local validation and scoring steps are linear in the size
of each generated scenario.

\subsection{System Architecture}

AdverSim is organized as a validate-and-requery generation loop followed by
offline evaluation layers. The generation layer accepts an attack type,
environment, and difficulty tuple, prompts the LLM for structured JSON, parses
the result, validates all ATT\&CK IDs against a fixed active matrix snapshot,
and re-queries when parsing, schema, or technique validation fails. Accepted
scenarios are written to the corpus and then passed to structural scoring,
SOC-simulation sensitivity analysis, Sigma rule generation and replay, real
baseline comparison, blinded human realism validation, and reproducibility CI.

\subsection{Algorithm 1: Validate-and-Requery Scenario Generation}

\begin{verbatim}
Input: parameter tuples X, active ATT&CK set A, max re-query R
Output: accepted corpus C and telemetry Q

C <- empty list
Q <- empty telemetry table

for each x = (attack_type, environment, difficulty) in X:
    prompt <- build_prompt(x)
    accepted <- false

    for attempt in 0..R:
        raw <- call_llm(prompt)
        parsed <- parse_json(raw)

        if parsed is invalid:
            prompt <- correction_prompt(x, "return valid JSON")
            record parse failure in Q
            continue

        issues <- validate_schema_and_attck(parsed, A)

        if issues is empty:
            C.append(parsed)
            record accepted scenario in Q
            accepted <- true
            break

        prompt <- correction_prompt(x, issues)
        record validation failure in Q

    if accepted is false:
        record rejection in Q

return C, Q
\end{verbatim}

The remote generation cost is
$O(N(R+1)\mathrm{LLM}(L))$. The local parsing and validation cost is
$O(N(R+1)(L+T+S+G))$, with local space $O(N(L+T+S+G)+K)$.

\subsection{Algorithm 2: Structural Realism Score}

For a scenario $s$, AdverSim computes active validity, structure, and transition
coherence:

\begin{verbatim}
tids <- ordered ATT&CK IDs from s.attack_stages
active_validity <- count(tid in A for tid in tids) / max(1, len(tids))

structure <- min(1,
    0.6 * len(stages) / 4
  + 0.4 * len(edges) / max(1, len(stages)-1))

transition_pairs <- adjacent unordered pairs from tids
transition_coherence <- count(pair in Pairs for pair in transition_pairs)
                        / max(1, len(transition_pairs))

R(S) <- 0.40 * active_validity
      + 0.35 * structure
      + 0.25 * transition_coherence
\end{verbatim}

Building the group-pair reference set costs $O(\sum_g k_g^2)$ for ATT\&CK
group technique counts $k_g$. Scoring one scenario costs $O(T+S+G)$ and uses
$O(P+T)$ space. The human validation experiment should be used to frame $R(S)$
as a structural consistency metric, not as a complete substitute for human
realism judgment.

\subsection{Algorithm 3: Human Realism Validation}

\begin{verbatim}
sample <- stratified_sample(C, n, attack_type x environment x difficulty)

for each scenario in sample:
    blind_id <- stable hash of seed and sample position
    write blinded scenario row
    write empty reviewer template row
    write hidden key row with R(S) features

collect H completed templates
aggregate ratings by blind_id
compute Krippendorff ordinal alpha by dimension
compute Spearman correlation between R(S) and mean overall realism
bootstrap the correlation confidence interval
fit exploratory R(S) weights on train split and evaluate held-out rows
\end{verbatim}

Packet generation costs $O(N+n(T+S+G))$. Rating aggregation costs $O(Hn)$,
agreement computation costs $O(nH^2)$, and bootstrap correlation costs
$O(Bn\log n)$.

\subsection{Algorithm 4: Detection Engineering and Replay}

\begin{verbatim}
for each selected scenario:
    focus_stage <- first high-signal stage, else first stage
    focus_tid <- first valid technique in focus_stage
    logsource <- map focus_tid to Sigma logsource
    keywords <- technique names and indicator terms
    write Sigma YAML rule

for each labelled event:
    labels <- extract ATT&CK labels
    haystack <- flattened lowercase event text
    for each Sigma rule:
        if logsource is compatible and keyword matches haystack:
            record alert and true-positive status
\end{verbatim}

Rule generation costs $O(M(T+L_s))$, where $L_s$ is stage text length. Replay
with transparent keyword matching costs $O(EMW)$ for average keyword count $W$.

\subsection{Complexity Summary}

\begin{table}[t]
\centering
\caption{AdverSim algorithmic complexity.}
\begin{tabular}{lll}
\hline
Component & Time & Space \\
\hline
Generation & $O(N(R+1)\mathrm{LLM}(L))$ remote & $O(N(L+T+S+G)+K)$ \\
Validation & $O(T+S)$ per scenario & $O(T+K)$ \\
$R(S)$ scoring & $O(\sum_g k_g^2 + N(T+S+G))$ & $O(P+T)$ \\
SOC sensitivity & $O(|\Theta|NS)$ & $O(|\Theta|D+N)$ \\
Sigma generation & $O(M(T+L_s))$ & $O(M)$ \\
Sigma replay & $O(EMW)$ & $O(M+E_{labels}+alerts)$ \\
Human validation & $O(N+n(T+S+G)+Hn+nH^2+Bn\log n)$ & $O(nH+P)$ \\
Figure CI & $O(A+F\cdot pixels)$ & $O(A+F\cdot pixels)$ \\
\hline
\end{tabular}
\end{table}
